\documentclass[12pt]{ximera}
\title{Practice Final Exam}
\begin{document}
\begin{abstract}
Practice for the cumulative final: which apportionment methods suffer which paradoxes
and quota violations, saddle points, a full four-method apportionment of teaching
aides, and expected value as a decision tool.
\end{abstract}
\maketitle

\section*{Practice Final Exam}

\begin{problem}
Assess whether each of the following claims is true or false.

\begin{prompt}
\textbf{(a)} Claim: it is possible for Hamilton's method of apportionment to apportion
more representation to a state than its upper quota.

\begin{multipleChoice}
\choice{True}
\choice[correct]{False}
\end{multipleChoice}

\begin{feedback}
False. Hamilton's method always satisfies the quota rule. Every state first receives its
lower quota, and the surplus seats hand out at most one extra seat apiece --- so each
state ends with either $\lfloor q\rfloor$ or $\lceil q\rceil$, never more.

Staying within quota is Hamilton's great virtue. Its weakness is the paradoxes, which
the divisor methods avoid.
\end{feedback}
\end{prompt}

\begin{prompt}
\textbf{(b)} Claim: the quotas of states with the largest populations change the fastest
when the standard divisor is modified.

\begin{multipleChoice}
\choice[correct]{True}
\choice{False}
\end{multipleChoice}

\begin{feedback}
True. A state's quota is $q = \frac{p}{d}$, so shrinking the divisor by a given
percentage raises every quota by that same percentage --- and a percentage of a large
quota is a larger absolute change.

For example, dropping the divisor from $26.13$ to $25.5$ raises a state of population
$1445$ from $55.30$ to $56.67$, a jump of $1.37$; the same change moves a state of
population $248$ only from $9.49$ to $9.73$, a jump of $0.24$.

This is why the divisor methods are biased: tuning the divisor moves the big states'
quotas past rounding boundaries much sooner than the small states'.
\end{feedback}
\end{prompt}

\begin{prompt}
\textbf{(c)} Claim: every game that can be represented with a $2\times 2$ payoff matrix
has a saddle point, where neither player has reason to deviate from that strategy.

\begin{multipleChoice}
\choice{True}
\choice[correct]{False}
\end{multipleChoice}

\begin{feedback}
False. The matrix
\[
\begin{array}{c|c|c|}
\phantom{|} & c & d \\\hline
a & 4 & -3 \\\hline
b & -1 & 1 \\\hline
\end{array}
\]
has maximin $-1$ and minimax $1$, which differ --- so there is no saddle point. The coin
game in Problem 4 is another example. Such games have no stable pure strategy, and both
players do better by mixing.
\end{feedback}
\end{prompt}

\begin{prompt}
\textbf{(d)} Claim: a zero-sum game is a game where one of the players has a winning
strategy, a way to guarantee themself victory regardless of the other player's moves.

\begin{multipleChoice}
\choice{True}
\choice[correct]{False}
\end{multipleChoice}

\begin{feedback}
False --- this confuses two different ideas.

\emph{Zero-sum} means only that the two players' payoffs always add to zero: whatever
one wins, the other loses. That is a statement about the payoff structure, not about
anybody's ability to force a win.

Plenty of zero-sum games give neither player a guaranteed win --- Rock/Paper/Scissors is
zero-sum and completely symmetric. The property being described in the claim belongs to
games like the Subtraction Game, where one player really can force a win.
\end{feedback}
\end{prompt}
\end{problem}

\begin{problem}
The following payoff matrices each represent Alice's perspective for a zero-sum game.
For each one, identify the saddle point.

\begin{prompt}
\textbf{(a)}
\[
\begin{array}{c|c|c|}
\phantom{|} & c & d \\\hline
a & -3 & 1 \\\hline
b & 0 & 2 \\\hline
\end{array}
\]
The saddle point value is $\answer{0}$.

\begin{feedback}
Row minima $-3$ and $0$ give a maximin of $0$; column maxima $0$ and $2$ give a minimax
of $0$. They agree, so the saddle point is $0$ at cell $(b,c)$.
\end{feedback}
\end{prompt}

\begin{prompt}
\textbf{(b)}
\[
\begin{array}{c|c|c|}
\phantom{|} & d & e \\\hline
a & 0 & -1 \\\hline
b & 1 & 2 \\\hline
c & -2 & 0 \\\hline
\end{array}
\]
The saddle point value is $\answer{1}$.

\begin{feedback}
Row minima are $-1$, $1$, $-2$, so the maximin is $1$. Column maxima are $1$ and $2$, so
the minimax is $1$. The saddle point is $1$ at cell $(b,d)$.

Note that row $b$ dominates both other rows for Alice, so she plays $b$; Bob then picks
$d$ to hold her to $1$ rather than $2$.
\end{feedback}
\end{prompt}

\begin{prompt}
\textbf{(c)}
\[
\begin{array}{c|c|c|c|}
\phantom{|} & d & e & f \\\hline
a & -1 & 5 & 0 \\\hline
b & -2 & 0 & -3 \\\hline
c & 2 & 3 & 1 \\\hline
\end{array}
\]
The saddle point value is $\answer{1}$.

\begin{feedback}
Row minima are $-1$, $-3$, and $1$, so the maximin is $1$. Column maxima are $2$, $5$,
and $1$, so the minimax is $1$. The saddle point is $1$ at cell $(c,f)$.

Here row $c$ dominates rows $a$ and $b$ for Alice, and column $f$ dominates $d$ and $e$
for Bob, so both players' choices are forced.
\end{feedback}
\end{prompt}
\end{problem}

\begin{problem}
Identify all applicable answers.

\begin{prompt}
\textbf{(a)} Which of the following can occur when using Hamilton's method of
apportionment?

\begin{selectAll}
\choice{Upper quota violation}
\choice[correct]{New States paradox}
\choice{Lower quota violation}
\choice[correct]{Population paradox}
\end{selectAll}

\begin{feedback}
The \textbf{New States paradox} and the \textbf{Population paradox} --- and the Alabama
paradox as well.

Hamilton's method never violates the quota rule, as Problem 1(a) established, so neither
quota violation is possible. What it does suffer is every one of the paradoxes, because
the surplus seats are handed out by comparing fractional parts, and those comparisons can
reshuffle when the house size, a population, or the set of states changes.
\end{feedback}
\end{prompt}

\begin{prompt}
\textbf{(b)} Which of the following can occur when using Adams' method of
apportionment?

\begin{selectAll}
\choice{Alabama paradox}
\choice{Upper quota violation}
\choice{New States paradox}
\choice[correct]{Lower quota violation}
\end{selectAll}

\begin{hint}
Adams' method rounds every modified quota \emph{up}. Which direction does that bias
push, and which states lose out?
\end{hint}

\begin{feedback}
Only the \textbf{lower quota violation}.

Adams' method is a divisor method, and divisor methods are immune to all three
paradoxes --- there is no fractional-parts comparison to reshuffle.

It can, however, break the quota rule, and always in the same direction. Adams rounds
up, which requires a divisor \emph{larger} than the standard one to avoid overshooting.
That shrinks every quota, and the large states --- whose quotas move fastest, by Problem
1(b) --- can be pushed below their lower quota. Homework 9 has an example: district A had
a standard quota of $86.915$ and Adams' method gave it only $85$.

It cannot exceed an upper quota: with a modified divisor $d \ge SD$ we get
$\frac{p}{d} \le q$, so $\lceil p/d\rceil \le \lceil q\rceil$. (Jefferson's method has
exactly the opposite bias --- it can break the upper quota but not the lower.)
\end{feedback}
\end{prompt}
\end{problem}

\begin{problem}
A school district needs to apportion $100$ teaching aides across four schools:

\begin{center}
\begin{tabular}{|r|c|c|c|c|c|}\hline
\textbf{School} & A & B & C & D & \textbf{Total}\\\hline
\textbf{Population} & 1{,}445 & 548 & 372 & 248 & 2{,}613 \\\hline
\end{tabular}
\end{center}

\begin{prompt}
\textbf{(a)} Compute the standard divisor and standard quotas, then complete the
apportionment with Hamilton's method.

Standard divisor: $\answer[tolerance=0.005]{26.13}$

\smallskip
$q_A = \answer[tolerance=0.0005]{55.3}$ \quad $q_B = \answer[tolerance=0.0005]{20.972}$
\quad $q_C = \answer[tolerance=0.0005]{14.237}$ \quad $q_D = \answer[tolerance=0.0005]{9.491}$

\smallskip
Hamilton's apportionment --- $A: \answer{55}$ \quad $B: \answer{21}$ \quad
$C: \answer{14}$ \quad $D: \answer{10}$

\begin{feedback}
\[ SD = \frac{2613}{100} = 26.13, \]
giving quotas $55.300$, $20.972$, $14.237$, and $9.491$, which total $100$. \checkmark

The lower quotas are $55, 20, 14, 9 = 98$, leaving $2$ surplus seats. The fractional
parts are
\[ A: .300, \qquad B: .972, \qquad C: .237, \qquad D: .491, \]
so B and D --- the two largest --- each gain a seat:
\[ A = 55, \quad B = 21, \quad C = 14, \quad D = 10. \]
\end{feedback}
\end{prompt}

\begin{prompt}
\textbf{(b)} Round the standard quotas as if you were doing just the first step of each
method. Give the resulting totals.

Jefferson (round down): $A: \answer{55}$, $B: \answer{20}$, $C: \answer{14}$,
$D: \answer{9}$, total $\answer{98}$

\smallskip
Adams (round up): $A: \answer{56}$, $B: \answer{21}$, $C: \answer{15}$, $D: \answer{10}$,
total $\answer{102}$

\smallskip
Webster (conventional): $A: \answer{55}$, $B: \answer{21}$, $C: \answer{14}$,
$D: \answer{9}$, total $\answer{99}$

\smallskip
Huntington-Hill cutoffs: $A: \answer[tolerance=0.005]{55.498}$,
$B: \answer[tolerance=0.005]{20.494}$, $C: \answer[tolerance=0.005]{14.491}$,
$D: \answer[tolerance=0.005]{9.487}$, giving a total of $\answer{100}$

\begin{feedback}
\[
\begin{array}{lccccl}
 & A & B & C & D & \text{total}\\
\text{quota} & 55.300 & 20.972 & 14.237 & 9.491 & 100\\
\text{Jefferson (down)} & 55 & 20 & 14 & 9 & 98\\
\text{Adams (up)} & 56 & 21 & 15 & 10 & 102\\
\text{Webster (round)} & 55 & 21 & 14 & 9 & 99
\end{array}
\]
Jefferson is $2$ short (needs a smaller divisor), Adams is $2$ over (needs a larger
one), and Webster is $1$ short.

For Huntington-Hill, the geometric means are
\[
\sqrt{55\cdot 56}=55.498, \quad \sqrt{20\cdot 21}=20.494, \quad
\sqrt{14\cdot 15}=14.491, \quad \sqrt{9\cdot 10}=9.487.
\]
Comparing: $55.300 < 55.498$ rounds down to $55$; $20.972 > 20.494$ rounds up to $21$;
$14.237 < 14.491$ rounds down to $14$; $9.491 > 9.487$ rounds up to $10$. That totals
exactly $100$, so Huntington-Hill needs no modified divisor at all --- and lands on the
same answer as Hamilton's method. (School D is a very close call: $9.491$ against a
cutoff of $9.487$.)
\end{feedback}
\end{prompt}

\begin{prompt}
\textbf{(c)} Complete the apportionment using each method, choosing the appropriate
modified divisor from $25.5$, $26.1$, and $26.7$.

\textbf{Jefferson's method} --- modified divisor $\answer[tolerance=0.05]{25.5}$, giving
$A: \answer{56}$, $B: \answer{21}$, $C: \answer{14}$, $D: \answer{9}$

\smallskip
\textbf{Adams' method} --- modified divisor $\answer[tolerance=0.05]{26.7}$, giving
$A: \answer{55}$, $B: \answer{21}$, $C: \answer{14}$, $D: \answer{10}$

\smallskip
\textbf{Webster's method} --- modified divisor $\answer[tolerance=0.05]{26.1}$, giving
$A: \answer{55}$, $B: \answer{21}$, $C: \answer{14}$, $D: \answer{10}$

\begin{hint}
Part (b) tells you which direction each divisor must move. Jefferson needs more seats,
so a smaller divisor; Adams needs fewer, so a larger one.
\end{hint}

\begin{feedback}
\textbf{Jefferson}, $d = 25.5$: modified quotas $56.667$, $21.490$, $14.588$, $9.725$;
rounding down gives $56, 21, 14, 9 = 100$. \checkmark

\textbf{Adams}, $d = 26.7$: modified quotas $54.120$, $20.524$, $13.933$, $9.288$;
rounding up gives $55, 21, 14, 10 = 100$. \checkmark

\textbf{Webster}, $d = 26.1$: modified quotas $55.364$, $20.996$, $14.253$, $9.502$;
conventional rounding gives $55, 21, 14, 10 = 100$. \checkmark

Comparing all five methods:
\[
\begin{array}{lcccc}
 & A & B & C & D\\
\text{Hamilton} & 55 & 21 & 14 & 10\\
\text{Jefferson} & 56 & 21 & 14 & 9\\
\text{Adams} & 55 & 21 & 14 & 10\\
\text{Webster} & 55 & 21 & 14 & 10\\
\text{Huntington-Hill} & 55 & 21 & 14 & 10
\end{array}
\]
Jefferson is the odd one out, moving a seat from the smallest school D to the largest
school A --- exactly the large-state bias you would expect from rounding everything down.
Every method here stays within quota.
\end{feedback}
\end{prompt}
\end{problem}

\begin{problem}
Use your knowledge of game theory and expected value to answer the following.

\begin{prompt}
\textbf{(a)} Complete the payoff matrix from Alice's perspective for a two-player
zero-sum game where Alice and Bob each have a quarter and a fifty-cent piece, and on the
count of three they each show one.
\begin{itemize}
\item If they show the same coin, Alice wins points equal to the sum of the coins.
\item If they show opposite coins, Bob wins points equal to the sum of the coins.
\end{itemize}

Give each entry in cents.
\[
\begin{array}{c|c|c|}
\phantom{|} & \text{Bob: }50\text{\textcent} & \text{Bob: }25\text{\textcent} \\\hline
\text{Alice: }50\text{\textcent} & \answer{100} & \answer{-75} \\\hline
\text{Alice: }25\text{\textcent} & \answer{-75} & \answer{50} \\\hline
\end{array}
\]

\begin{feedback}
\begin{itemize}
\item Both $50$\textcent: same, Alice wins $100$\textcent\ $\to +100$.
\item Alice $50$\textcent, Bob $25$\textcent: opposite, Bob wins $75$\textcent\ $\to -75$.
\item Alice $25$\textcent, Bob $50$\textcent: opposite $\to -75$.
\item Both $25$\textcent: same, Alice wins $50$\textcent\ $\to +50$.
\end{itemize}
This game has no saddle point --- the maximin is $-75$ and the minimax is $50$ --- which
is the counterexample for Problem 1(c).
\end{feedback}
\end{prompt}

\begin{prompt}
\textbf{(b)} Alice is deciding whether to go to a basketball game that her friend Jalen
might be playing in. There is a $30\%$ chance he doesn't get to play, a $30\%$ chance he
plays but they lose, and a $40\%$ chance he plays and they win.

\[
\begin{array}{c|c|c|c|}
\phantom{|} & \text{doesn't play} & \text{plays}+\text{lose} & \text{plays}+\text{win} \\\hline
\text{go to the game} & -10 & 2 & 7 \\\hline
\text{stay home} & 10 & -5 & -3 \\\hline
\end{array}
\]

Expected value of going: $\answer[tolerance=0.005]{0.4}$ \qquad
Expected value of staying home: $\answer[tolerance=0.005]{0.3}$

\smallskip
Which should she choose?
\begin{multipleChoice}
\choice[correct]{Go to the game}
\choice{Stay home}
\end{multipleChoice}

\begin{feedback}
\begin{align*}
E(\text{go}) &= 0.30(-10) + 0.30(2) + 0.40(7) = -3.0 + 0.6 + 2.8 = 0.4,\\
E(\text{stay home}) &= 0.30(10) + 0.30(-5) + 0.40(-3) = 3.0 - 1.5 - 1.2 = 0.3.
\end{align*}
Since $0.4 > 0.3$, Alice should go to the game --- though by the narrowest of margins.
A small revision to any of her ratings would flip the recommendation, which is worth
remembering: expected value always produces a definite answer, but not always a robust
one.
\end{feedback}
\end{prompt}
\end{problem}

\end{document}
